\par\medskip\noindent\textbf{Cómo leer la salida.}
{\small\setlength{\tabcolsep}{4pt}\renewcommand{\arraystretch}{1.15}
\begin{xltabular}{\linewidth}{@{}>{\raggedright\arraybackslash}p{0.24\linewidth}>{\raggedleft\arraybackslash}p{0.15\linewidth}>{\raggedright\arraybackslash}X@{}}
\toprule Indicador & Valor & Interpretación \\ \midrule \endfirsthead
\toprule Indicador & Valor & Interpretación \\ \midrule \endhead
\bottomrule \endlastfoot
Combinación (media) & $R^2_{OS}=0{,}61$\,\%; CW $p=0{,}017$ & Promediar los diez pronósticos supera a la media histórica de forma significativa. \\
Combinación (mediana) & 0,50\,\%; $p=0{,}011$ & Resultado similar. \\
Componentes principales & $-0{,}36$ a $-1{,}53$\,\% & Resumir los predictores en componentes no mejora el ECM, aunque CW es significativo (pronósticos correlacionados con el objetivo pero con varianza excesiva). \\
Desv. est. de pronósticos & 0,110 / 0,082 & La combinación se mueve poco alrededor de la media: es su principal virtud. \\
Correlación con el realizado & 0,073 & Baja, pero suficiente para mejorar el pronóstico. \\
Equivalente cierto, media histórica & 6,62\,\% & Rendimiento seguro que valora igual el inversionista con aversión 3. \\
Combinación & 7,57\,\%; ganancia 0,95 & Usar la combinación vale casi un punto por año: valor económico relevante aunque el $R^2$ sea pequeño. \\
Componentes ($k=1$) & ganancia 0,39 & Valor económico positivo pero menor. \\
Solo dp & ganancia $-0{,}88$ & Seguir solo a la razón dividendo/precio habría costado casi un punto anual. \\
Peso medio en acciones & 1,17 / 0,94 / 0,52 / 0,37 & Exposición promedio de cada estrategia. \\
\end{xltabular}}
